% !TEX root = ./main.tex
\chapter*{Conclusion}
\addcontentsline{toc}{chapter}{Conclusion}

Ce mémoire s'est donné pour objectif de construire un cadre cohérent de
modélisation et d'estimation du risque de défaut à l'aide du modêle
\gls{cir}, puis d'en évaluer la pertinence pour la valorisation
des Credit Default Swaps à partir de données de marché réelles.

Dans un premier temps, nous avons posé les fondations théoriques
nécessaires à l'étude : présentation des \gls{cds} et du risque de crédit,
rappels de calcul stochastique et de statistique. Nous avons ensuite
développé, dans un espace probabilisé filtré, un cadre rigoureux de
modélisation du défaut fondé sur la notion d'intensité, établi la
formule de réduction permettant de ramener le calcul de flux contingents au
défaut à une espérance sous la mesure risque-neutre, puis justifié le
choix du modêle \gls{cir} au regard des propriétés de positivité, de retour à
la moyenne et de tractabilité analytique qu'il requiert. Ce cadre nous a
permis de dériver les expressions de la jambe fixe et de la jambe de
protection d'un \gls{cds}, et donc de son spread théorique, ainsi qu'une méthode
d'estimation des paramêtres du modêle par la méthode des moments appliquée
à la distribution stationnaire du processus.

L'application numérique menée sur trois entreprises présentant des profils
de risque contrastés, Johnson \& Johnson, Motorola Solutions et Bombardier
Inc, a permis de valider l'implémentation du modêle : les probabilités de
survie calculées analytiquement coïncident avec celles obtenues par
simulation de Monte-Carlo, avec des écarts négligeables. L'analyse de
sensibilité a par ailleurs confirmé le rôle économique attendu de chacun
des paramêtres du modêle : le niveau moyen de long terme $b$ gouverne le
niveau général des spreads longs, la vitesse de retour à la moyenne $a$
en rêgle la dynamique de convergence, l'intensité initiale $\lambda_0$
domine le court terme, et le taux de recouvrement $\delta$ produit un effet
d'échelle quasi proportionnel sur l'ensemble de la courbe.

Ces résultats confirment que le modêle \gls{cir} constitue un outil pertinent et
économiquement interprétable pour une premiêre modélisation du risque de
crédit. Il présente cependant des limites significatives, qui sont
apparues clairement dans notre étude. D'une part, le modêle monofactoriel
retenu peine à reproduire la structure par termes observée sur le
marché : les courbes de spreads simulées restent relativement plates,
tandis que les courbes observées, en particulier pour les entités les plus
risquées comme Bombardier, présentent une pente nettement marquée, ce dont
témoignent les erreurs quadratiques moyennes obtenues. D'autre part, pour
certaines entités, les paramêtres calibrés ne satisfont pas la condition
de Feller, ce qui se traduit par des trajectoires simulées visitant
fréquemment le voisinage de zéro avant de connaître des variations
brutales, un comportement mathématiquement admissible mais moins convaincant
d'un point de vue économique. Plus généralement, le caractêre continu du
processus \gls{cir}, l'hypothêse de paramêtres constants dans le temps et
l'hypothêse d'indépendance entre taux d'intérêt et intensité de défaut
constituent autant de limites intrinsêques au modêle.

Plusieurs pistes permettraient de prolonger ce travail. Sur le plan de
l'estimation, le recours à la méthode du maximum de vraisemblance ou à des
approches bayésiennes permettrait d'exploiter davantage la dynamique
temporelle des observations, plutot que de s'appuyer sur la seule
distribution stationnaire du processus. Sur le plan de la modélisation,
l'introduction de paramêtres dépendant du temps ou de composantes à sauts,
comme dans les modêles de type Jump-\gls{cir}, offrirait une plus grande
flexibilité pour représenter les changements de régime et les variations
brutales de spread. Enfin, le recours à des techniques d'apprentissage
automatique, en complément ou en substitution des approches paramétriques
classiques, constitue une voie de recherche prometteuse pour améliorer la
calibration et la valorisation des dérivés de crédit.

En définitive, ce mémoire montre que le modêle \gls{cir}, par son équilibre
entre simplicité, interprétabilité économique et tractabilité
analytique, demeure un point de départ pertinent pour la modélisation du
risque de crédit, tout en soulignant la nécessité de développer des
modêles plus riches pour rendre pleinement compte de la complexité des
marchés du crédit contemporains.